\documentclass[12pt]{article}
\usepackage[utf8]{inputenc}
\usepackage{amsmath, amssymb, amsthm}
\usepackage{logicproof}
\usepackage[most]{tcolorbox}
\usepackage{xcolor}
\usepackage{geometry}

\geometry{a4paper, margin=0.8in}


\title{Sample Test 1: Logical Proof}
\author{\textbf{Nguyen Dinh Thien Loc} \\ Student ID: 24125093}

\begin{document}
\maketitle
\section*{Problem 2}
\subsection*{(a)}
\begin{logicproof}{2}
\exists x (P(x) \wedge Q(x)) & premise \\
\forall x (Q(x) \rightarrow R(x)) & premise \\
\begin{subproof}
    x_0 \quad P(x_0) \wedge Q(x_0) & assumption \\
    Q(x_0) \rightarrow R(x_0) & $\forall x e$ 2 \\
    Q(x_0) & $\wedge e_2$ 3 \\
    R(x_0) & $\rightarrow e$ 4, 5 \\
    P(x_0) & $\wedge e_1$ 3 \\
    P(x_0) \wedge R(x_0) & $\wedge i$ 7, 6 \\
    \exists x (P(x) \wedge R(x)) & $\exists x i$ 8
\end{subproof}
\exists x (P(x) \wedge R(x)) & $\exists x e$ 1, 3--9
\end{logicproof}

\subsection*{(b)}
\begin{logicproof}{3}
\exists x P(x) & premise \\
\forall x \forall y (P(x) \rightarrow Q(y)) & premise \\
\begin{subproof}
    y_0 & \\
    \begin{subproof}
        x_0 \quad P(x_0) & assumption \\
        \forall y (P(x_0) \rightarrow Q(y)) & $\forall x e$ 2 \\
        P(x_0) \rightarrow Q(y_0) & $\forall y e$ 5 \\
        Q(y_0) & $\rightarrow e$ 6, 4
    \end{subproof}
    Q(y_0) & $\exists x e$ 1, 4--7
\end{subproof}
\forall y Q(y) & $\forall y i$ 3--8
\end{logicproof}

\section*{Problem 3}

\subsection*{(a) $\forall x(P(x) \rightarrow Q(f(x)))$}

\begin{tcolorbox}[title=Model $\mathcal{M}_1$ (Satisfies the formula)]
    \begin{itemize}
        \item \textbf{Domain} $\mathcal{D} = \{1, 2\}$
        \item \textbf{Function} $f(n) = n$ (Identity function)
        \item \textbf{Predicates:} $P = \{1\}$, $Q = \{1, 2\}$
        \item \textbf{Reason:} For $x=1$, $P(1) \rightarrow Q(1)$ is $T \rightarrow T$ (True). For $x=2$, $P(2) \rightarrow Q(2)$ is $F \rightarrow T$ (True).
    \end{itemize}
\end{tcolorbox}

\begin{tcolorbox}[title=Model $\mathcal{M}_2$ (Does not satisfy)]
    \begin{itemize}
        \item \textbf{Domain} $\mathcal{D} = \{1\}$
        \item \textbf{Function} $f(1) = 1$
        \item \textbf{Predicates:} $P = \{1\}$, $Q = \emptyset$
        \item \textbf{Reason:} For $x=1$, $P(1) \rightarrow Q(1)$ becomes $T \rightarrow F$, which is False.
    \end{itemize}
\end{tcolorbox}

\subsection*{(b) $\exists x \forall y ((P(x,y) \wedge P(y,x)) \vee (Q(x,y) \wedge R(y,x)))$}

\begin{tcolorbox}[title=Model $\mathcal{M}_3$ (Satisfies the formula)]
    \begin{itemize}
        \item \textbf{Domain} $\mathcal{D} = \{a\}$
        \item \textbf{Predicates:} $P = \{(a,a)\}$, $Q = \emptyset$, $R = \emptyset$
        \item \textbf{Reason:} If we pick $x=a$, then for all $y$ (which can only be $a$), the first part of the disjunction $(P(a,a) \wedge P(a,a))$ is True. Thus the whole formula is True.
    \end{itemize}
\end{tcolorbox}

\begin{tcolorbox}[title=Model $\mathcal{M}_4$ (Does not satisfy)]
    \begin{itemize}
        \item \textbf{Domain} $\mathcal{D} = \{1, 2\}$
        \item \textbf{Predicates:} $P = \emptyset$, $Q = \emptyset$, $R = \emptyset$
        \item \textbf{Reason:} All predicates are empty relations. For any $x$ we choose, both $(P \wedge P)$ and $(Q \wedge R)$ will be False for all $y$, making the entire disjunction False.
    \end{itemize}
\end{tcolorbox}

\newpage
\section*{Problem 4}

\subsection*{(a) $\exists x P(x) \wedge \exists x Q(x) \models \exists x (P(x) \wedge Q(x))$}
\textbf{Result:} This semantic entailment \textbf{does not hold}.

\begin{tcolorbox}[title=Counter-example Model]
    \begin{itemize}
        \item \textbf{Domain} $\mathcal{D} = \{1, 2\}$
        \item \textbf{Predicates:} $P = \{1\}$, $Q = \{2\}$
        \item \textbf{Reason:} 
            \begin{itemize}
                \item $\exists x P(x)$ is True (witness $x=1$).
                \item $\exists x Q(x)$ is True (witness $x=2$).
                \item Thus, the premise $\exists x P(x) \wedge \exists x Q(x)$ is True.
                \item However, there is no element in $\mathcal{D}$ that is in both $P$ and $Q$. Thus, $\exists x (P(x) \wedge Q(x))$ is False.
            \end{itemize}
    \end{itemize}
\end{tcolorbox}

\subsection*{(b) $\exists y \forall x P(x,y) \models \forall x \exists y P(x,y)$}
\textbf{Result:} This semantic entailment \textbf{holds}.

\begin{tcolorbox}[title=Proof Sketch]
    Assume the premise $\exists y \forall x P(x,y)$ is true in a model $\mathcal{M}$ with environment $l$. This means there exists some element $c \in \mathcal{D}$ such that for all elements $a \in \mathcal{D}$, the pair $(a, c) \in P^{\mathcal{M}}$. 
    
    To show $\forall x \exists y P(x,y)$, we must show that for any arbitrary $a \in \mathcal{D}$, there exists some $b \in \mathcal{D}$ such that $(a, b) \in P^{\mathcal{M}}$. For any such $a$, we can choose $b = c$. Since the premise guarantees $(a, c) \in P^{\mathcal{M}}$ for all $a$, the conclusion is satisfied.
\end{tcolorbox}

\newpage
\section*{Problem 5}

\subsection*{(a) $\phi = ((p \rightarrow q) \wedge p) \vee q$}

To check for validity, we apply the Tableau Algorithm to the negation $\neg\phi$:
$$\neg (((p \rightarrow q) \wedge p) \vee q)$$

\begin{tcolorbox}[title=Tableau Path for $\neg\phi$]
\begin{enumerate}
    \item $C_0 = \{ \neg (((p \rightarrow q) \wedge p) \vee q) \}$
    \item $C_1 = \{ \neg ((p \rightarrow q) \wedge p), \neg q \}$ \hfill ($\neg\vee$-rule)
    \item \textbf{Path Branching} ($\neg\wedge$-rule applied to $\neg ((p \rightarrow q) \wedge p)$):
    \begin{itemize}
        \item \textbf{Branch 1:} $C_2 = \{ \neg (p \rightarrow q), \neg q \}$
            \begin{itemize}
                \item $C_3 = \{ p, \neg q, \neg q \}$ \hfill ($\neg\rightarrow$-rule)
                \item This path is \textbf{complete} and \textbf{not closed} (no clash).
            \end{itemize}
        \item \textbf{Branch 2:} $C_2 = \{ \neg p, \neg q \}$
            \begin{itemize}
                \item This path is \textbf{complete} and \textbf{not closed} (no clash).
            \end{itemize}
    \end{itemize}
\end{enumerate}
\end{tcolorbox}

\textbf{Conclusion:} Since there exists a complete tableau path that is not closed, the negated formula $\neg\phi$ is \textbf{satisfiable}. Therefore, the original formula $\phi$ is \textbf{not valid}.

\subsection*{(b) $\psi = \forall x \neg P(x) \vee \exists x( P(x) \vee Q(x))$}

We apply the Tableau Algorithm to the negation $\neg\psi$:
$$\neg (\forall x \neg P(x) \vee \exists x( P(x) \vee Q(x)))$$

\begin{tcolorbox}[title=Tableau Path for $\neg\psi$]
\begin{enumerate}
    \item $C_0 = \{ \neg (\forall x \neg P(x) \vee \exists x( P(x) \vee Q(x))) \}$
    \item $C_1 = \{ \neg \forall x \neg P(x), \neg \exists x (P(x) \vee Q(x)) \}$ \hfill ($\neg\vee$-rule)
    \item $C_2 = \{ \exists x \neg \neg P(x), \neg \exists x (P(x) \vee Q(x)) \}$ \hfill ($\neg\forall$-rule)
    \item $C_3 = \{ \exists x P(x), \neg \exists x (P(x) \vee Q(x)) \}$ \hfill ($\neg\neg$-rule)
    \item $C_4 = \{ \exists x P(x), \forall x \neg (P(x) \vee Q(x)) \}$ \hfill ($\neg\exists$-rule)
    \item $C_5 = \{ P(c), \forall x \neg (P(x) \vee Q(x)) \}$ \hfill ($\exists$-rule, creates new constant $c$)
    \item $C_6 = \{ P(c), \neg (P(c) \vee Q(c)) \}$ \hfill ($\forall$-rule using $c$)
    \item $C_7 = \{ P(c), \neg P(c), \neg Q(c) \}$ \hfill ($\neg\vee$-rule)
\end{enumerate}
\end{tcolorbox}

\textbf{Conclusion:} The path $C_7$ contains a **clash** $\{P(c), \neg P(c)\}$. Since all possible fair tableau paths are closed, the negated formula $\neg\psi$ is \textbf{unsatisfiable}. Therefore, the original formula $\psi$ is \textbf{valid} (a tautology).

\end{document}
